% 一次线性化更新：训练输入梯度产生参数增量，另一输入梯度读取该增量。
% 两个梯度向量只示意方向与内积，不表示测量值。
\begin{tikzpicture}[foundation picture]
 \path[use as bounding box] (0,2) rectangle (112,69);
 \draw[foundation axis] (2,10)--(42,10);
 \draw[foundation axis] (2,10)--(2,54);
 \draw[foundation flow,draw=FoundationInput] (2,10)--(33,43);
 \draw[foundation flow,draw=FoundationModel!80!black] (2,10)--(38,20);
 \node at (22,49) {$\bm\varphi_0(\bm x_i)$};
 \node at (28,13) {$\bm\varphi_0(\bm x_j)$};
 \node[foundation loss,minimum width=30mm,minimum height=8mm] (loss) at (80,62) {损失导数$a_i$};
 \node[foundation model,text width=52mm,minimum height=12mm] (update) at (80,43)
  {$\Delta\bm v=-\eta a_i\bm\varphi_0(\bm x_i)$};
 \draw[foundation flow] (loss.south)--(update.north);
 \draw[foundation flow] (33,43)--(update.west);
 \node[foundation output,text width=52mm,minimum height=16mm] (response) at (80,20)
  {$\Delta f(\bm x_j)=\langle\bm\varphi_0(\bm x_j),\Delta\bm v\rangle$\\
   $=-\eta a_iK_0(\bm x_j,\bm x_i)$};
 \draw[foundation flow] (update.south)--(response.north);
 \draw[foundation flow] (38,20)--(response.west);
\end{tikzpicture}%
